\documentclass[11pt,a4paper]{article}
\usepackage[utf8]{inputenc}
\usepackage[T1]{fontenc}
\usepackage{lmodern}
\usepackage[margin=2.4cm]{geometry}
\usepackage{graphicx}
\usepackage{booktabs}
\usepackage{amsmath}
\usepackage{siunitx}
\usepackage[table]{xcolor}
\usepackage[colorlinks=true,linkcolor=black,urlcolor=blue!60!black,citecolor=black]{hyperref}
\usepackage{caption}
\captionsetup{font=small,labelfont=bf}



\graphicspath{{../figures/}}
\sisetup{detect-all}
\newcommand{\ns}{\,\mathrm{ns}}

\title{\bfseries The $\gamma$-flash time base of the n\_TOF EAR2 X17 campaign\\
\large A calibration from the seven runs taken with the SiPM-wall divert disabled}
\author{n\_TOF X17 / DREAM analysis}
\date{2026-07-28}

\begin{document}
\maketitle

\begin{abstract}
\noindent
The SiPM walls of the X17 setup are blanked around the $\gamma$-flash, so they
never record it, and the pulse-shape analysis (PSA) that produces the official
hit files mis-times the flash on every detector except the beam pickup. We
identify the seven runs of the 2026 EAR2 campaign in which the blanking gate was
inadvertently disabled --- the only direct view of the true flash arrival in the
walls --- and use them to calibrate the flash time of all 32 wall channels
against the beam pickup~(PKUP):
\[
   t_\mathrm{flash}(\text{bunch}) \;=\; t_\mathrm{PKUP}(\text{bunch}) + C ,
   \qquad C \simeq -1719\ns .
\]
The constant reproduces to $0.5\ns$ between runs of one epoch, and the liquid
scintillators --- which are never gated --- show that the underlying time base
is stable to $0.68\ns$ rms over the whole campaign. We further show that the
$5\ns$ shift between parasitic and dedicated proton pulses is an artefact of
front-end saturation rather than a change of the flash arrival time, and we
quantify why the plastic scintillators cannot be used for flash timing.
\end{abstract}

\section{Result}

The calibration is a per-channel constant $C$ referred to the pickup pulse of
the same bunch:
\begin{equation}
   t_\mathrm{flash\ at\ detector}(\text{bunch}) = t_\mathrm{PKUP}(\text{bunch}) + C[\text{channel}] .
   \label{eq:cal}
\end{equation}
Table~\ref{tab:walls} gives the per-wall values; the per-channel numbers are in
\texttt{data/flash\_timing\_calibration.json}. The recommended set is the
2026-07-16 epoch, for runs $\geq 224400$.

\begin{table}[htbp]
\centering
\small
\caption{Wall constants $C = t_\mathrm{channel}-t_\mathrm{PKUP}$, in ns.
``spread'' is the r.m.s.\ over the eight channels of a wall.}
\label{tab:walls}
\begin{tabular}{lrrrrrr}
\toprule
wall & all 7 runs & 2026-07-11 & \textbf{2026-07-16} & channel spread & per-bunch $\sigma$ & flash amp.\ [ADC] \\
\midrule
WALA & $-1717.5$ & $-1716.7$ & $\mathbf{-1719.5}$ & 5.2 & 3.2 & 28613 \\
WALB & $-1719.3$ & $-1719.1$ & $\mathbf{-1719.7}$ & 3.2 & 3.1 & 28515 \\
WALC & $-1715.5$ & $-1714.0$ & $\mathbf{-1719.2}$ & 1.9 & 3.6 & 28295 \\
WALD & $-1714.1$ & $-1712.2$ & $\mathbf{-1719.0}$ & 3.9 & 3.4 & 28246 \\
\midrule
all 32 & $-1716.6$ & $-1715.5$ & $\mathbf{-1719.3}$ & 3.4 & 3.5 & 28420 \\
\bottomrule
\end{tabular}
\end{table}

\begin{table}[htbp]
\centering
\small
\caption{Uncertainty budget.}
\label{tab:budget}
\begin{tabular}{lr}
\toprule
effect & size [ns] \\
\midrule
run-to-run reproducibility within an epoch & 0.5 \\
time-base stability over the campaign (LIQ common mode, r.m.s.) & 0.68 \\
transport from the 07-16 epoch to the end of the campaign & 1.2 \\
$07\text{-}11 \rightarrow 07\text{-}16$ epoch difference (real hardware change) & 3.9 \\
beam-intensity dependence, parasitic $\rightarrow$ dedicated & 5.0 \\
\midrule
single-bunch spread, one channel & 3.5 \\
single-bunch spread, 32-channel average & 2.5 \\
\bottomrule
\end{tabular}
\end{table}

\paragraph{Do not use the stored \texttt{tflash}.} In gated runs the PSA flash
finder is bistable \emph{within a single run}: in some bunches it tags the
blanking-gate transient, in others the clamped leak at the true flash time. No
per-run offset can repair that. The pickup finder, by contrast, never fails
(\S\ref{sec:pkup}).

\section{Why special runs were needed}

The wall signal is diverted around the flash by a ${\sim}1\,\mu$s blanking gate
(N1081B module~6, section~C). The walls therefore record either the gate
transient --- $1190\ns$ before the flash in the July-11 configuration, $375\ns$
after the gate delay was changed from 200 to $1000\ns$ on 2026-07-22 --- or a
clamped leak at the true flash time.

A scan of every processed run of the campaign (306 runs, 224264--224587)
identifies exactly seven runs in which the gate was inactive
(Table~\ref{tab:runs}): all four walls then tag the same time to within $\pm6\ns$
and the flash peak reaches ${\sim}28\,500$~ADC instead of the ${\sim}850$~ADC
transient. The 07-11 group corresponds to the switch outage during which the
N1081B module~6 was offline; the 07-16 pair coincides with the wall/plastic
recalibration. 145 runs before 224345 have no wall trees in the official
processing and cannot be checked without a tape reprocessing.

\begin{table}[htbp]
\centering
\small
\caption{The seven divert-off runs.}
\label{tab:runs}
\begin{tabular}{lllrl}
\toprule
run & date & bunches & \multicolumn{1}{c}{intensity} & all-four-walls agree \\
\midrule
224356 & 07-10 21:44 -- 07-11 01:00 & 3412 & dedicated & 85\,\% \\
\textbf{224357} & 07-11 01:01 -- 03:51 & 3286 & 47\,\% ded.\ / 53\,\% par. & 98\,\% \\
224358 & 07-11 03:52 -- 06:40 & 3269 & mixed & 93\,\% \\
224359 & 07-11 06:41 -- 09:40 & 3426 & dedicated & 87\,\% \\
224360 & 07-11 09:41 -- 10:08 & 549 & dedicated & 91\,\% \\
224464 & 07-16 09:54 -- 12:56 & 2716 & mostly dedicated & 60\,\% \\
224466 & 07-16 ${\sim}$13:00 -- 15:19 & 2073 & mostly dedicated & 100\,\% \\
\bottomrule
\end{tabular}
\end{table}

\section{Method}

For every (bunch, channel) we take the largest-amplitude hit within $\pm300\ns$
of that bunch's flash anchor and subtract the pickup pulse time of the same
bunch. The quoted constant is the mean over bunches after a $\pm100\ns$ core cut
that removes PSA mis-tags. Since \texttt{tof} is quantised to 1\,ns, the mean
over ${\sim}3000$ bunches has a statistical error of $0.06\ns$: every number
below is limited by systematics, not statistics.

\section{Per-channel constants and resolution}

\begin{figure}[htbp]
\centering
\includegraphics[width=\textwidth]{01_per_channel_arrival.png}
\caption{Per-channel flash arrival referred to the pickup, for all seven runs
(top), and the decomposition of the run-to-run scatter into within-epoch
reproducibility and the real 07-11$\rightarrow$07-16 hardware shift (bottom).
Within an epoch every channel reproduces to ${\sim}0.6\ns$; the epoch shift is
channel-dependent and reaches $9\ns$ on WALD.}
\label{fig:perchan}
\end{figure}

On 07-11 the four walls sat $7\ns$ apart; by 07-16 they agree to $0.7\ns$. The
shift is per-channel (WALB $-0.6\ns$, WALD $-6.9\ns$), i.e.\ a front-end change
rather than a cable change, consistent with the wall threshold/HV equalisation
of 07-15/16. The equalised state is the one that held for the rest of the
campaign, which is why the 07-16 constants are recommended.

\begin{figure}[htbp]
\centering
\includegraphics[width=\textwidth]{04_distributions.png}
\caption{Left: single channel against the pickup. Right: the difference of two
channels of the same wall, in which the pickup and beam jitter cancel.}
\label{fig:dist}
\end{figure}

Decomposing the 32-channel matrix bunch by bunch gives, averaged over the seven
runs: $3.5\ns$ for a single channel against the pickup, of which $2.5\ns$ is
common to all 32 channels and $2.5\ns$ is independent. The independent part
averages away over the array; the common part does not, and since the pickup
itself contributes only ${\sim}1\ns$ (\S\ref{sec:pkup}) it represents a genuine
bunch-to-bunch variation of the flash arrival. The 07-16 runs are better than
the 07-11 ones ($2.7$ vs $3.9\ns$ per channel), again following the
equalisation.

\begin{figure}[htbp]
\centering
\includegraphics[width=\textwidth]{03_stability.png}
\caption{Within-run stability: pooling all 32 channels in 40 bins across a full
three-hour run, the bin-to-bin spread is $0.25$--$0.39\ns$ and the maximum
excursion $1.4\ns$.}
\label{fig:stab}
\end{figure}

\subsection{The reference: the beam pickup}
\label{sec:pkup}

The pickup is a near-ideal reference: its pulse time has a bunch-to-bunch spread
of $0.83\ns$ (dedicated) to $1.42\ns$ (parasitic), it is usable in 98.4\,\% of
bunches --- the 1.6\,\% of failures are all parasitic pulses --- its amplitude
is strictly linear in protons (ratio $1.97$ for a factor two), and its timing
moves by only $-0.45\ns$ between the two intensity classes.

\section{Is the intensity dependence real?}

The wall constant differs by $-5.0\ns$ between parasitic ($4.1\times10^{12}$~p)
and dedicated ($8.5\times10^{12}$~p) pulses. Three independent arguments show
that this is a property of the measurement, not a change of the flash arrival
time.

\begin{figure}[htbp]
\centering
\includegraphics[width=\textwidth]{06_intensity_mechanism.png}
\caption{(a) The shift is strongly detector-dependent --- $0$ for the pickup,
$-5\ns$ for the walls, $-3$ to $-6\ns$ for the plastics, $-16\ns$ for the
silicon monitor. A genuine change of the flash arrival would be common to all.
(b) Only the pickup is linear in protons; every other detector is saturated, its
amplitude ratio being $\simeq 1$. (c) The shift follows the change in pulse
risetime.}
\label{fig:mech}
\end{figure}

\begin{enumerate}
\item \textbf{It is detector-dependent.} The same flash is seen by all
detectors, so a real change in its arrival time would shift them equally.
Instead the shift ranges from $0$ (pickup) through $-5\ns$ (walls) to $-16\ns$
(silicon monitor, which is the most deeply saturated, with the PSA saturation
flag set on 78\,\% of its flash hits).

\item \textbf{The pickup does not move.} The pickup timestamps the very protons
that produce the flash, and its amplitude doubles between the two classes while
its time moves by $-0.45\ns$. The flash cannot precede its own protons.

\item \textbf{Classical amplitude walk is excluded quantitatively.} Within a
single intensity class the dependence of the hit time on its amplitude is
measurable; extrapolating it to the observed change of median amplitude between
classes predicts $+0.3\ns$ for the walls and ${<}0.4\ns$ for the others --- with
the wrong sign for the walls (Fig.~\ref{fig:pred}). The measured shifts are ten
to fifty times larger.
\end{enumerate}

\begin{figure}[htbp]
\centering
\includegraphics[width=0.62\textwidth]{07_intensity_prediction.png}
\caption{The amplitude dependence measured \emph{within} one intensity class
cannot account for the shift \emph{between} classes.}
\label{fig:pred}
\end{figure}

The mechanism is saturation of the front end. The peak amplitude cannot grow
with the flash size, so a larger flash instead reaches the clipping level
earlier and stays there longer: on the walls the reported risetime falls from
$20.5$ to $15.4\ns$ ($-25\,\%$) while the FWHM grows from ${\sim}700$ to
${\sim}2000\ns$; on the silicon monitor the risetime falls from $98$ to
$56\ns$. Any edge-based estimator therefore fires earlier. \emph{Practical
consequence:} the effect must be carried in the calibration (use the constant
matching the pulse type, or apply the $-5\ns$ per class), but it must not be
interpreted as physics.

\section{Stability in time}

Because the walls can only be measured in the seven divert-off runs, the
stability of the time base is established with the liquid scintillators and the
plastics, which are never gated and therefore see the flash in every run.

\begin{figure}[htbp]
\centering
\includegraphics[width=\textwidth]{11_timeseries.png}
\caption{Flash time relative to the pickup, run by run, with statistical errors;
bottom panel: the single-bunch spread. The liquids are stable, the plastics are
not.}
\label{fig:ts}
\end{figure}

Over 36--39 runs spanning 224400--224584 (2026-07-13 to 07-28) the four liquid
cells give
\[
 \sigma_\mathrm{run}(\mathrm{LIQA}) = 0.73\ns,\quad
 \sigma_\mathrm{run}(\mathrm{LIQB}) = 0.73\ns,\quad
 \sigma_\mathrm{run}(\mathrm{LIQC}) = 0.77\ns,\quad
 \sigma_\mathrm{run}(\mathrm{LIQD}) = 1.48\ns .
\]
The four cells move \emph{together}: the within-run spread of their deviations
is $0.22\ns$ and the LIQA--LIQC correlation across runs is $+0.92$, so what is
observed is a common wander of the time base rather than detector noise. Its
r.m.s.\ is $0.68\ns$ and the largest single-run excursion is $-2.2\ns$
(run~224531). The two divert-off runs of 07-16 sit $+0.9$ to $+1.2\ns$ from the
late-campaign mean, which is the figure quoted in Table~\ref{tab:budget} for the
transport of the wall calibration to the end of the campaign.

\paragraph{Independent check in gated runs.} Where the gate is active the walls
still emit a clamped leak. Measured against Eq.~\eqref{eq:cal} in runs 224362
(07-11 era) and 224572 (07-26, fifteen days after the nearest calibration run
and after the gate-delay change), the leak lands within $6.2\ns$ on average and
$9.2\ns$ worst case on all 32 channels. The gate transient itself sits at
$-1190\ns$ and $-375\ns$ respectively, reproducing the $+800\ns$ delay change of
07-22.

\section{Why the plastics cannot be used --- in the official processing}

\begin{quote}\small
\textbf{Note added.} Everything in this section describes the \emph{official}
processed files. The parallel UserInput work (\texttt{FINDINGS\_2026-07-28\_psa\_optimization.md})
reprocessed run 224572, and on the \texttt{v4\_walshapes} variant the plastic
flash hit is found in \textbf{100\,\%} of bunches with $\sigma = 3.0$--$4.6\ns$,
against 42\,\% and $2.5\ns$ officially. The single change responsible is the
\textbf{G-FLASH threshold} ($50 \rightarrow 2000/10^4$): the \texttt{v1\_flash}
variant, with the elimination window untouched, already reaches 100\,\%. The
missing hits were therefore a downstream symptom of the flash-finder bug --- with
the flash mis-located at ${\sim}130\ns$ (fitting noise), the pulse recognition
for that bunch never reconstructs the real pulse at $11.63\,\mu$s. Widening the
elimination window (\texttt{v2\_elim}) changes the yield not at all, and the
surviving flash hits have area/amp $\simeq 72$--$76$, above even the widened
window --- so elimination is not the gate.
\end{quote}

\begin{figure}[htbp]
\centering
\includegraphics[width=\textwidth]{08_plastic_anatomy.png}
\caption{Hit density per bunch in the flash region. The wall produces one clean
pulse; the plastic response is spread over ${\sim}100\ns$ and accompanied by a
continuum of fragments. In the post-FIFO run the wall shows both the gate
transient (${\sim}-2100\ns$) and the leak (${\sim}-1750\ns$).}
\label{fig:anat}
\end{figure}

The plastics fail in three distinct ways, all quantified in
\texttt{data/plastic\_pathology.csv}:

\begin{table}[htbp]
\centering
\small
\caption{Flash-region behaviour by detector family and epoch. ``found'' is the
fraction of bunches in which the PSA put \emph{any} hit within $60\ns$ of the
modal flash position; $\sigma$ and the pulse quantities refer to the
largest-amplitude hit there; ``$n_\mathrm{hit}$'' counts hits per bunch in the
$\pm1\,\mu$s flash window.}
\label{tab:path}
\begin{tabular}{llrrrrrr}
\toprule
epoch & family & found & $\sigma$ [ns] & amplitude & sat.\ flag & FWHM [ns] & $n_\mathrm{hit}$/bunch \\
\midrule
2026-07-11 & WAL & 100\,\% & 4.1 & 28521 & 0\,\% & 1527 & 12.4 \\
2026-07-11 & PSS & \textbf{67\,\%} & \textbf{19.1} & 23193 & 0\,\% & 11 & 38.2 \\
2026-07-16 & WAL & 100\,\% & 2.3 & 28560 & 0\,\% & 2602 & 12.5 \\
2026-07-16 & PSS & 64\,\% & 5.8 & 34571 & 0\,\% & 43 & 26.0 \\
2026-07-16 & LIQ & \textbf{100\,\%} & 4.1 & \textbf{63931} & \textbf{80\,\%} & 103 & 43.6 \\
2026-07-26 & WAL & 100\,\% & 8.3 & 19650 & 6\,\% & 57 & 11.7 \\
2026-07-26 & PSS & \textbf{42\,\%} & 2.5 & 55114 & 27\,\% & 76 & 24.3 \\
2026-07-26 & LIQ & \textbf{100\,\%} & 4.5 & \textbf{63921} & \textbf{79\,\%} & 102 & 44.1 \\
\bottomrule
\end{tabular}
\end{table}

\paragraph{(i) Yield.} A hit at the flash time is reconstructed in only
64--67\,\% of bunches, falling to \textbf{42\,\%} after the FIFO fan-in was
installed on 2026-07-17, when the plastic flash pulse begins to rail the
digitiser (median amplitude 55\,100, up to 63\,600 of a 65\,536 range). In the
bunches where it is missing, the PSA reconstructs \emph{no hit at all} before
${\sim}-1600\ns$; we verified that this is not an artefact of the extraction
window, which covered the flash region in 100\,\% of those bunches.

\paragraph{The liquids are the control --- and they exclude saturation as the
cause.} LIQ sits on the same digitiser and the same PSA, and its flash hits are
\emph{more} saturated than the plastics' (saturation flag on 79--80\,\% against
27\,\%, amplitude 63\,900 against 55\,100), with the highest hit multiplicity of
all (44 per bunch). Yet the PSA finds the liquid flash in \textbf{100\,\%} of
bunches at $\sigma = 4.1$--$4.5\ns$. The plastic drop-out is therefore specific
to the plastic channel --- its PSA configuration and/or its pulse shape --- and
not a generic consequence of a railed pulse, which suggests it can be repaired
in the \texttt{UserInput}.

\paragraph{(ii) Amplitude-dependent bias and spread.} When a hit is found, its
time depends on its amplitude: in run 224357 the median moves by $17\ns$ between
the lowest and highest amplitude bins, with a per-bunch $\sigma$ of $7$--$15\ns$
(Fig.~\ref{fig:pamp}). After the FIFO the precision improves markedly
($\sigma = 2.4$--$3.4\ns$, bias $4\ns$) but the yield collapses, so the two
regimes trade against each other.

\begin{figure}[htbp]
\centering
\includegraphics[width=\textwidth]{09_plastic_vs_amp.png}
\caption{Plastic timing precision (left) and bias (right) versus the amplitude
of the flash hit, for the three epochs.}
\label{fig:pamp}
\end{figure}

\paragraph{(iii) Instability in time.} The plastic constant moves by tens of
nanoseconds over the campaign ($\sigma_\mathrm{run} = 17$--$31\ns$,
Fig.~\ref{fig:ts}) as the plastic HV was equalised (07-16) and the FIFO fan-in
installed (07-17). Channel PSSB\,ch2 is defective throughout (amplitude 9\,778
against ${\sim}30\,000$ elsewhere). Plastic constants must therefore be taken
per run, never per epoch.

\begin{figure}[htbp]
\centering
\includegraphics[width=\textwidth]{10_plastic_channels.png}
\caption{Per-channel flash-hit yield and timing spread for the eight plastic
readouts in the three reference runs.}
\label{fig:pchan}
\end{figure}

\section{Independent cross-check against the reprocessed data}

The reprocessing provides a second, independent route to the same constants: a
\emph{corrected} flash finder, run on 224572 --- a gated run, ten days after the
nearest calibration run, in which the walls can only see the clamped leak.
Table~\ref{tab:xval} compares its stored \texttt{tflash}, referenced to the
pickup (8 partials, 1597 bunches), against this work.

\begin{table}[htbp]
\centering
\small
\caption{Reprocessed (\texttt{v4\_walshapes}) stored \texttt{tflash} minus the
pickup, against the constants of this work. All values in ns.}
\label{tab:xval}
\begin{tabular}{lrrrl}
\toprule
tree & reprocessed & this work & difference & this work is from \\
\midrule
WALA & $-1717.82$ & $-1719.46$ & $+1.64$ & divert-off runs, 07-16 epoch \\
WALB & $-1718.07$ & $-1719.69$ & $+1.62$ & divert-off runs, 07-16 epoch \\
WALC & $-1718.05$ & $-1719.23$ & $+1.18$ & divert-off runs, 07-16 epoch \\
WALD & $-1722.29$ & $-1719.01$ & $-3.28$ & divert-off runs, 07-16 epoch \\
\multicolumn{3}{r}{mean $+0.29$, rms} & \textbf{2.07} & \\
\midrule
LIQA & $-1708.38$ & $-1708.22$ & $-0.16$ & campaign monitor (official proc.) \\
LIQB & $-1710.78$ & $-1710.30$ & $-0.48$ & campaign monitor (official proc.) \\
LIQC & $-1695.73$ & $-1695.56$ & $-0.17$ & campaign monitor (official proc.) \\
LIQD & $-1701.27$ & $-1701.56$ & $+0.29$ & campaign monitor (official proc.) \\
\multicolumn{3}{r}{mean $-0.13$, rms} & \textbf{0.27} & \\
\bottomrule
\end{tabular}
\end{table}

The walls agree to $2.1\ns$ rms and the liquids to $0.27\ns$ rms --- the latter
being two different processings of the same run agreeing to a quarter of a
nanosecond. Both sit inside the transport budget of Table~\ref{tab:budget}, and
the two routes share no assumption beyond the pickup itself.

\section{How to apply, and what limits it}

\begin{enumerate}
\item \textbf{Walls}: use the per-channel \texttt{C\_2026\_07\_16} for runs
$\geq 224400$, \texttt{C\_2026\_07\_11} before that. Add $-5\ns$ if the pulse
type differs from the one the constant was taken on.
\item \textbf{Liquids}: the per-run value, or the campaign mean; they agree to
${\sim}1\ns$.
\item \textbf{Plastics}: per-run value only.
\item Never use the stored \texttt{tflash} of WAL/PSS/LIQ \emph{from the
official files}. In reprocessed files it is sound, but PKUP\,$+\,C$ remains the
recommended time base: it is measured on the undiverted flash, whereas even a
reprocessed wall \texttt{tflash} times the clamped leak of a gated signal.
\item If a bunch has no usable pickup pulse (1.6\,\%, all parasitic), drop the
bunch rather than falling back on \texttt{tflash}.
\end{enumerate}

\paragraph{Limitations.} (i) The undiverted flash saturates the wall front end,
so these runs give timing only --- the energy scale must still come from the
$^{88}$Y runs 224476--224479. (ii) The $3.9\ns$ shift between 07-11 and 07-16 is
real and per-channel, so wall front-end timing can move; nothing guarantees that
it did not move again between 07-16 and 07-28 beyond the liquid-scintillator
evidence that the shared time base did not. (iii) Runs before 224345 are
unverified. (iv) The raw data of all seven runs has aged off EOS disk; a
waveform-level study needs a tape stage.

\end{document}
